\chapter{Particle Filters}
\label{ch:bf-particle-filters}

Particle filters approximate the filtering distribution with weighted
particles rather than a Gaussian moment closure.  They are important for
nonlinear and non-Gaussian state-space models, but they occupy a different
place in BayesFilter's HMC program than EKF or sigma-point filters.  Standard
resampling is discrete, and differentiating a particle-filter likelihood is
not the same problem as differentiating a Kalman recursion.

\section{Separation of Roles}
\label{sec:bf-pf-separation}

BayesFilter should separate three roles:
\begin{enumerate}[label=(\roman*)]
  \item a particle filter as a likelihood estimator for likelihood comparison
    or particle MCMC;
  \item a differentiable particle filter as an approximate differentiable
    filtering backend;
  \item a particle method as a simulation-based diagnostic for nonlinear
    filters.
\end{enumerate}
These roles have different validity conditions.  A stochastic likelihood
estimator may be useful for pseudo-marginal methods, but BayesFilter does not
yet have accepted primary-source support in this monograph pass for strong
pseudo-marginal HMC claims.  A differentiable resampling relaxation may support
gradients, but it generally changes the filtering approximation and needs its
own bias and convergence assumptions.

\section{Differentiable Frontier}
\label{sec:bf-differentiable-pf-frontier}

The Phase 2B literature gate accepted the entropy-regularized optimal-transport
particle-filter source only as frontier support.  It justifies discussing
differentiable particle filters under their stated assumptions.  It does not
certify arbitrary differentiable resampling, and it does not make particle
filters a drop-in HMC backend for BayesFilter.

Therefore this chapter should remain conservative until the literature gate is
completed:
\begin{itemize}
  \item do not claim generic unbiased differentiable likelihood gradients;
  \item do not claim pseudo-marginal HMC correctness without primary PMCMC and
    pseudo-marginal support;
  \item do not replace the linear Gaussian derivative spine with a particle
    approximation when exact or controlled approximate filters are available;
  \item use particle methods as reference diagnostics only when the Monte Carlo
    error is measured and reported.
\end{itemize}

\section{Validation Requirements}
\label{sec:bf-pf-validation-requirements}

A BayesFilter particle backend needs diagnostics beyond sampler summaries:
effective sample size, log-likelihood estimator variance, resampling counts,
random-seed reproducibility policy, particle degeneracy warnings, and
comparison to exact or dense-reference cases where possible.  If a relaxed
resampling method is used, the validation report must state the relaxed target
being differentiated and how it differs from the intended filtering target.
